\documentclass[10pt,a4paper]{amsart}
\usepackage{amsmath,amssymb,amsthm,graphicx}
% Retained minimal preamble for the Stacks Project chapter flat.tex (More on Flatness).
% Verbatim excerpt of the chapter preamble: Xy-pic setup (native preamble lines 14-18) and the native
% theorem declarations with the mathematical macros (native preamble lines 154-217). Copied mechanically;
% nothing added, removed or redefined. The wrapper class and the amsmath/amssymb/amsthm/graphicx base
% packages are supplied by the assembler, so the class/externaldocument part of the native preamble is omitted.
% For commutative diagrams we use Xy-pic
\usepackage[all]{xy}

% We use 2cell for 2-commutative diagrams.
\xyoption{2cell}
% Theorem environments.
%
\theoremstyle{plain}
\newtheorem{theorem}[subsection]{Theorem}
\newtheorem{proposition}[subsection]{Proposition}
\newtheorem{lemma}[subsection]{Lemma}

\theoremstyle{definition}
\newtheorem{definition}[subsection]{Definition}
\newtheorem{example}[subsection]{Example}
\newtheorem{exercise}[subsection]{Exercise}
\newtheorem{situation}[subsection]{Situation}

\theoremstyle{remark}
\newtheorem{remark}[subsection]{Remark}
\newtheorem{remarks}[subsection]{Remarks}

\numberwithin{equation}{subsection}

% Macros
%
\def\lim{\mathop{\mathrm{lim}}\nolimits}
\def\colim{\mathop{\mathrm{colim}}\nolimits}
\def\Spec{\mathop{\mathrm{Spec}}}
\def\Hom{\mathop{\mathrm{Hom}}\nolimits}
\def\Ext{\mathop{\mathrm{Ext}}\nolimits}
\def\SheafHom{\mathop{\mathcal{H}\!\mathit{om}}\nolimits}
\def\SheafExt{\mathop{\mathcal{E}\!\mathit{xt}}\nolimits}
\def\Sch{\mathit{Sch}}
\def\Mor{\mathop{\mathrm{Mor}}\nolimits}
\def\Ob{\mathop{\mathrm{Ob}}\nolimits}
\def\Sh{\mathop{\mathit{Sh}}\nolimits}
\def\NL{\mathop{N\!L}\nolimits}
\def\CH{\mathop{\mathrm{CH}}\nolimits}
\def\proetale{{pro\text{-}\acute{e}tale}}
\def\etale{{\acute{e}tale}}
\def\QCoh{\mathit{QCoh}}
\def\Ker{\mathop{\mathrm{Ker}}}
\def\Im{\mathop{\mathrm{Im}}}
\def\Coker{\mathop{\mathrm{Coker}}}
\def\Coim{\mathop{\mathrm{Coim}}}

% Boxtimes
%
\DeclareMathSymbol{\boxtimes}{\mathbin}{AMSa}{"02}

%
% Macros for moduli stacks/spaces
%
\def\QCohstack{\mathcal{QC}\!\mathit{oh}}
\def\Cohstack{\mathcal{C}\!\mathit{oh}}
\def\Spacesstack{\mathcal{S}\!\mathit{paces}}
\def\Quotfunctor{\mathrm{Quot}}
\def\Hilbfunctor{\mathrm{Hilb}}
\def\Curvesstack{\mathcal{C}\!\mathit{urves}}
\def\Polarizedstack{\mathcal{P}\!\mathit{olarized}}
\def\Complexesstack{\mathcal{C}\!\mathit{omplexes}}
% \Pic is the operator that assigns to X its picard group, usage \Pic(X)
% \Picardstack_{X/B} denotes the Picard stack of X over B
% \Picardfunctor_{X/B} denotes the Picard functor of X over B
\def\Pic{\mathop{\mathrm{Pic}}\nolimits}
\def\Picardstack{\mathcal{P}\!\mathit{ic}}
\def\Picardfunctor{\mathrm{Pic}}
\def\Deformationcategory{\mathcal{D}\!\mathit{ef}}
\usepackage{cleveref}
\crefname{originallemma}{original lemma}{original lemmas}
\Crefname{originallemma}{Original lemma}{Original lemmas}
\makeatletter
\newlabel{lemma-base-change-complete}{{05HJ}{source}{}{sourcepointer-1}{}}
\newlabel{lemma-base-change-complete@cref}{{[originallemma][0][]05HJ}{[1][0][]source}{}{}{}}
\makeatother
\title{Existence of a complete devissage after one elementary etale localization}
\author{The Stacks Project Authors}
\date{stacks-project:flat.tex:local-snapshot-20261011; GNU FDL 1.2 or later}
\begin{document}
\maketitle
\noindent\textit{Editorial note.} All delivered mathematical text is the author's own, copied from the pinned Stacks Project chapter file flat.tex; no statement or proof step was rewritten, reordered or removed.\par\medskip
\noindent\textit{Editorial note.} The preamble is the chapter's own retained minimal preamble (sources/stacks-flat/retained-preamble.tex), a verbatim excerpt of the native Xy-pic setup and of the native theorem declarations and mathematical macros; the Xy-pic xy environment is required because the counted proof displays two commutative diagrams. The native class line, the xr-hyper externaldocument block and the reference/slogan/history display environments are omitted because the assembler supplies a plain amsart wrapper and this excerpt uses none of them; no macro was invented or redefined.\par\medskip
\noindent\textit{Editorial note.} The notion of a complete devissage is defined earlier in the same chapter and is declared as an allowed standard prerequisite instead of being dragged in as a context span, because the counted obligation is the existence and gluing argument, not the definition.\par\medskip
\noindent\textit{Allowed standard prerequisites (editorial).}\begin{itemize}
\item Complete devissage of a finite type module (paper's own notion). Same chapter, Definition environment labelled definition-complete-devissage, source line 1195; the pinned file defines it and it is quoted in the statement as the property being produced. Supplies the meaning of 'there exists a complete devissage of F\_i/Y\_i/S' over s'' and the recognised data (Z\_k, Y\_k, i\_k, pi\_k, G\_k, alpha\_k); no part of the counted construction is taken from it.
\item Elementary etale neighbourhood of a pointed scheme. Earlier chapters of the same project; imported by this chapter as standard input (Stacks tags 057K, 057L, 057R are the concrete predecessors used). Supplies the notion of the pointed morphism (S', s') -> (S, s) and the fact that its underlying morphism g : X' -> X is open, used at proof lines 1596-1601 to cover the fibre.
\end{itemize}\medskip
\ifdefined\hypertarget\hypertarget{sourcepointer-1}{}\fi
\noindent\textit{Source reference.} Original lemma 05HJ: sources/stacks-flat/pinned-original-source.tex, source lines 1271--1271. Same-chapter navigation pointer to the base-change lemma invoked at proof line 1614 for the finitely many Y\_i; the counted content of the proof spans (pointwise data, finite cover of X\_s, fibre product of the bases) is retained verbatim in the delivered body, so this pointer replaces nothing.\par\medskip
\begin{lemma}
\label{lemma-existence-complete}
Let $X \to S$ be a finite type morphism of schemes.
Let $\mathcal{F}$ be a finite type quasi-coherent $\mathcal{O}_X$-module.
Let $s \in S$ be a point.
There exists an elementary \'etale neighbourhood
$(S', s') \to (S, s)$ and \'etale morphisms
$h_i : Y_i \to X_{S'}$, $i = 1, \ldots, n$ such that for each
$i$ there exists a complete d\'evissage of $\mathcal{F}_i/Y_i/S'$ over $s'$,
where $\mathcal{F}_i$ is the pullback of $\mathcal{F}$ to $Y_i$
and such that $X_s = (X_{S'})_{s'} \subset \bigcup h_i(Y_i)$.
\end{lemma}

\begin{proof}
For every point $x \in X_s$ we can find a diagram
$$
\xymatrix{
(X, x) \ar[d] & (X', x') \ar[l]^g \ar[d] \\
(S, s) & (S', s') \ar[l]
}
$$
of pointed schemes such that the horizontal arrows are elementary
\'etale neighbourhoods and such that $g^*\mathcal{F}/X'/S'$ has a
complete d\'evissage at $x'$. As $X \to S$ is of finite type the
fibre $X_s$ is quasi-compact, and since each $g : X' \to X$ as above
is open we can cover $X_s$ by a finite union of $g(X'_{s'})$.
Thus we can find a finite family of such diagrams
$$
\vcenter{
\xymatrix{
(X, x) \ar[d] & (X'_i, x'_i) \ar[l]^{g_i} \ar[d] \\
(S, s) & (S'_i, s'_i) \ar[l]
}
}
\quad i = 1, \ldots, n
$$
such that $X_s = \bigcup g_i(X'_i)$. Set
$S' = S'_1 \times_S \ldots \times_S S'_n$
and let $Y_i = X_i \times_{S'_i} S'$ be the base change of $X'_i$ to $S'$. By
Lemma \ref{lemma-base-change-complete}
we see that the pullback of $\mathcal{F}$ to $Y_i$ has a complete d\'evissage
over $s$ and we win.
\end{proof}

\end{document}
